\par\textbf{Empirical Strategy}\par

\par\textbf{Twitter: Panel OLS and IV with Fixed Effects}\par

\begin{itemize}
\item
  The baseline specification (equation 1) is:\footnote{(Binstok, Gould, and Kaplan 2024, p. 2337, l. 4--12)}
\end{itemize}

\emph{Israel Negative\textsubscript{jt} = β\textsubscript{0} + β\textsubscript{1} Jews Negative\textsubscript{jt} + μ\textsubscript{j} + δ\textsubscript{t} + ε\textsubscript{jt}}

\begin{itemize}
\item
  μ\textsubscript{j} denotes district fixed effects (355 districts); δ\textsubscript{t} denotes year, month-of-year and day-of-week fixed effects.\footnote{(Binstok, Gould, and Kaplan 2024, p. 2337, l. 10--12)}
\item
  β\textsubscript{1} measures the effect of the local negative salience of Jews on negative chatter about Israel.\footnote{(Binstok, Gould, and Kaplan 2024, p. 2337, l. 13--15)}
\end{itemize}

\begin{itemize}
\item
  The unit of observation is the district-day rather than the individual: only 2.5\% of users posting about Israel also post about Jews on the same day, so tweets serve as a proxy for local sentiment.\footnote{(Binstok, Gould, and Kaplan 2024, p. 2337, l. 15--25)}
\item
  Identification rests on changes within districts over time, after district and calendar effects have been removed.
\item
  In the IV specification, negative tweets about Tottenham instrument for negative tweets about Jews.\footnote{(Binstok, Gould, and Kaplan 2024, p. 2337, l. 25--26; p. 2340, l. 2--6)}

  \begin{itemize}
  \item
    \textbf{Relevance:} negative tweets about Tottenham predict negative tweets about Jews (Table 2, columns 6--7)\footnote{(Binstok, Gould, and Kaplan 2024, p. 2339, l. 2--8)}, with a first-stage F-statistic of 25.83.\footnote{(Binstok, Gould, and Kaplan 2024, p. 2340, l. 10--12; Table 3, column 1)}
  \item
    \textbf{Support for exclusion:} chatter about Tottenham does not predict chatter about China (placebo).\footnote{(Binstok, Gould, and Kaplan 2024, p. 2339, l. 14--24)}
  \item
    Chatter about Tottenham follows the timing and location of the club's matches.\footnote{(Binstok, Gould, and Kaplan 2024, p. 2339, l. 25--35)}
  \item
    Alternative instruments are dummies for match windows and the locations of opposing teams.\footnote{(Binstok, Gould, and Kaplan 2024, p. 2340, l. 19--29)}
  \end{itemize}
\item
  Standard errors are robust and clustered at the district level.\footnote{(Binstok, Gould, and Kaplan 2024, notes to Tables 2 and 3)}
\item
  Because of the large number of zeros, the model is re-estimated at the weekly and monthly level and with Poisson fixed effects.\footnote{(Binstok, Gould, and Kaplan 2024, p. 2342, l. 27--31; p. 2343, l. 1--7 and l. 17--22)}
\end{itemize}

\par\textbf{ALLBUS: Cross-Sectional OLS}\par

\begin{itemize}
\item
  The outcome is regressed on the variable of interest, with controls added step by step (Table 7, columns 1--7, p. 2353).\footnote{(Binstok, Gould, and Kaplan 2024, p. 2352, l. 11--23)}
\item
  The authors argue that the coefficient remains stable or rises as controls are added, which makes it unlikely that omitted variables bias it away from zero.\footnote{(Binstok, Gould, and Kaplan 2024, p. 2352, l. 34--37; p. 2354, l. 27--32)}

  \begin{itemize}
  \item
    The same pattern holds within age, gender and education subgroups.\footnote{(Binstok, Gould, and Kaplan 2024, p. 2355, l. 4--13)}
  \item
    Bounds following Oster (2019) and Krauth (2016) suggest that the true effect is larger.\footnote{(Binstok, Gould, and Kaplan 2024, p. 2355, l. 14--19)}
  \end{itemize}
\item
  The authors name three remaining threats: further omitted variables, reverse causality and measurement error, as respondents may conceal antisemitic views.\footnote{(Binstok, Gould, and Kaplan 2024, p. 2355, l. 21--28)}
\end{itemize}

\par\textbf{ALLBUS: Instrumental Variable Design}\par

\begin{itemize}
\item
  The structural equation underlying Table 10 can be written as:
\end{itemize}

\emph{Israel Unjust\textsubscript{i} = β Jews Influence\textsubscript{i} + δ Born1930--60\textsubscript{i} + γ′W\textsubscript{i} + ε\textsubscript{i}}

\begin{itemize}
\item
  The corresponding first stage is:
\end{itemize}

\emph{Jews Influence\textsubscript{i} = π Born1930--45\textsubscript{i} + δ Born1930--60\textsubscript{i} + γ′W\textsubscript{i} + u\textsubscript{i}}

\begin{itemize}
\item
  The underlying idea is that children who grew up under Nazi rule (1933--1945) were indoctrinated against Jews and that this effect persisted throughout their lives.\footnote{(Binstok, Gould, and Kaplan 2024, pp. 2355--2356, l. 29--36 and l. 2--3)}

  \begin{itemize}
  \item
    Figure 7 shows that respondents born in the 1930s and 1940s are more likely to agree that Jews have too much influence.\footnote{(Binstok, Gould, and Kaplan 2024, p. 2356, l. 4--7)}
  \item
    As the State of Israel did not exist at that time, the indoctrination was directed against Jews and not against Israel.\footnote{(Binstok, Gould, and Kaplan 2024, p. 2356, l. 11--12)}
  \end{itemize}
\item
  Identification compares respondents born 1930--1945 with those born 1946--1960, holding the broad cohort (before or after 1960) constant.\footnote{(Binstok, Gould, and Kaplan 2024, p. 2358, l. 1--6)}
\item
  The endogenous variable is ``Jews have too much influence''; the instrument is the indicator for being born 1930--1945.\footnote{(Binstok, Gould, and Kaplan 2024, p. 2358, l. 34--37)}
\item
  Controls are added column by column (Table 10, columns 1--8, p. 2359), and standard errors are robust.\footnote{(Binstok, Gould, and Kaplan 2024, p. 2360, l. 7--16)}
\end{itemize}

\par\textbf{Identifying Assumptions}\par

\begin{itemize}
\item
  \textbf{Relevance:} the instrument must shift antisemitic views.

  \begin{itemize}
  \item
    With a linear age control, the 1930--1945 cohort is significant but, in the authors' words, ``not strong enough''.\footnote{(Binstok, Gould, and Kaplan 2024, p. 2356, l. 7--9 and l. 17--20)}
  \item
    The first stage becomes considerably stronger when the indicator for being born before 1960 is used instead; this version is used for the main IV (Table 9, third panel: 0.519, standard error 0.118).\footnote{(Binstok, Gould, and Kaplan 2024, p. 2356, l. 20--21; p. 2358, l. 1--7)}
  \end{itemize}
\item
  \textbf{Exogeneity and exclusion restriction:} being born 1930--1945 rather than 1946--1960 affects views on Israel only through antisemitism.\footnote{(Binstok, Gould, and Kaplan 2024, p. 2358, l. 10--15)}

  \begin{itemize}
  \item
    The authors argue that the year of birth is ``exogenously determined''.\footnote{(Binstok, Gould, and Kaplan 2024, p. 2356, l. 15--16)}
  \item
    They acknowledge that cohorts differ in politics, religion, nationalism and attitudes towards Islam and foreigners, which may affect views on Israel directly.\footnote{(Binstok, Gould, and Kaplan 2024, p. 2358, l. 19--31)}
  \end{itemize}
\item
  \textbf{Linearity and homogeneous effects:} the model implies one coefficient per step on the seven-point scale. The paper does not discuss this assumption explicitly, so it is a point for the evaluation.
\item
  \textbf{Twitter IV:} chatter about Tottenham affects chatter about Israel only through the salience of Jews, and Tottenham is not associated with Israel.\footnote{(Binstok, Gould, and Kaplan 2024, p. 2333, l. 1--8; p. 2339, l. 20--24)}
\end{itemize}

\par\textbf{The Authors' Defence of the Assumptions}\par

\begin{itemize}
\item
  The indicator for being born before 1960 allows for a broad direct effect of age; its coefficient is insignificant.\footnote{(Binstok, Gould, and Kaplan 2024, p. 2358, l. 12--15; p. 2360, l. 1--6)}
\item
  The more restrictive specification without any age control (column 1) yields the smallest estimate, so a violation of this assumption would, if anything, bias the estimate towards zero.\footnote{(Binstok, Gould, and Kaplan 2024, p. 2360, l. 1--3; fn. 23, pp. 2360--2361, l. 42--44 and l. 41--43)}
\item
  Adding worldview controls that are correlated with age increases the estimate.\footnote{(Binstok, Gould, and Kaplan 2024, p. 2360, l. 17--33)}

  \begin{itemize}
  \item
    The argument follows the logic of coefficient stability used by Alesina, Giuliano and Nunn (2013) and Autor, Dorn and Hanson (2013).\footnote{(Binstok, Gould, and Kaplan 2024, p. 2360, l. 30--32)}
  \end{itemize}
\item
  Excluding respondents born 1940--1945 barely changes the result (0.90 instead of 0.85).\footnote{(Binstok, Gould, and Kaplan 2024, fn. 21, p. 2360, l. 36--39)}
\item
  Alternative instrument 1: year of birth (linear), estimated separately within the cohorts born before and after 1960 (Table 12).\footnote{(Binstok, Gould, and Kaplan 2024, p. 2361, l. 33--40; p. 2363, l. 33--41)}
\item
  Alternative instrument 2: an interviewer aged 75 or older, with whom respondents are assumed to be more open, yields similar results, but the first-stage F-statistic is only about 5.\footnote{(Binstok, Gould, and Kaplan 2024, pp. 2364--2365, l. 36--44 and l. 1--6)}
\end{itemize}

